% Expository adaptation, October 7, 2026. Original mathematics: OpenAI, September 23, 2026.
% Reorganized and expanded for pedagogy; see README.md for the pinned source and scope.
\section{Historical context and related work}


The complexity of probabilistic Turing machines was developed by
Gill~\cite{Gill1977}, including measures of both time and tape usage.
Random walks made the logarithmic-space question concrete:
Aleliunas, Karp, Lipton, Lov\'asz and Rackoff obtained polynomial-time
randomized algorithms for undirected connectivity and explicitly asked
whether polynomial-time randomized logarithmic space could be simulated
in deterministic logarithmic space~\cite{AleliunasKLLR1979}.
Early work also studied probabilistic machines without a polynomial-time
restriction. Borodin, Cook and Pippenger obtained a quadratic-space
deterministic simulation in this broader setting
\cite{BorodinCookPippenger1983}.
The distinction between a space bound and simultaneous time and space
bounds is essential throughout this history.

Pseudorandom generators attack the problem by replacing a long sequence
of fresh random bits with the output generated from a short seed.
Ajtai, Koml\'os and Szemer\'edi gave early logarithmic-space
constructions that hit small-space tests using short random
strings~\cite{AjtaiKomlosSzemeredi1987}. Babai, Nisan and Szegedy
subsequently obtained pseudorandom generators through multiparty
communication bounds~\cite{BabaiNisanSzegedy1992}.
Nisan's generator uses $O(S\log R)$ seed bits for a space-$S$
computation reading $R$ random bits~\cite{Nisan1992}.
For logarithmic space and polynomially many bits this gives
$O(\log^2 n)$ seed length; direct seed enumeration has a
quasipolynomial time bound. Nisan's later simultaneous simulation
achieves polynomial time and $O(\log^2 n)$ space
\cite{Nisan1994RLinSC}. Although its title uses $\mathrm{RL}$,
that result includes two-sided bounded error.
For $S(n)\ge\log n$, Nisan and Zuckerman showed that a space-$S$
computation using $\operatorname{poly}(S)$ random bits can be simulated
using only $O(S)$ random bits and $O(S)$ space~\cite{NisanZuckerman1996}.
The restriction on the original number of random bits is substantive;
for logarithmic space it permits polylogarithmically many bits.

Saks and Zhou reduced the deterministic space needed for general
bounded-error logarithmic-space computation to
$O(\log^{3/2}n)$~\cite{SaksZhou1999}.
Cai, Chakaravarthy and van Melkebeek established a simultaneous
tradeoff: for each fixed rational $0\le\alpha\le1/2$, the simulation
uses time $n^{O(\log^{1/2-\alpha}n)}$ and space
$O(\log^{3/2+\alpha}n)$
\cite{CaiChakaravarthyVanMelkebeek2006}.
The polynomial-time endpoint has squared-logarithmic space.
Hoza subsequently improved the general space bound to
$O(\log^{3/2}n/\sqrt{\log\log n})$
\cite[Theorem~1.6 and Corollary~A.5]{Hoza2021}.
These smaller space containments do not assert the same polynomial-time
endpoint. Hoza's work also develops weighted pseudorandomness, in which
signed or real weights replace an ordinary probability distribution;
the space-containment improvement and the new weighted-generator
construction are distinct results of that paper.

Acceptance-probability estimation also connects derandomization to
stochastic matrix products and inverses. Ahmadinejad, Kelner, Murtagh,
Peebles, Sidford and Vadhan developed high-precision small-space
algorithms for random-walk probabilities, with stronger bounds for
Eulerian graphs~\cite{AhmadinejadKelnerMurtaghPeeblesSidfordVadhan2020}.
Cohen, Doron, Sberlo and Ta-Shma later used varying precision across
recursion to improve the space complexity of long stochastic matrix
products~\cite{CohenDoronSberloTaShma2023}. The construction in these notes works with a finite inverse on a time-layered configuration graph;
its correction and copy hierarchy controls the information retained
while individual entries are evaluated.

Structural restrictions on a computation provide another route.
For undirected connectivity, Nisan, Szemer\'edi and Wigderson obtained
deterministic space $O(\log^{3/2}n)$~\cite{NisanSzemerediWigderson1992},
subsequently improved to $O(\log^{4/3}n)$ by Armoni, Ta-Shma,
Wigderson and Zhou~\cite{ArmoniTaShmaWigdersonZhou2000}.
Reingold's logarithmic-space algorithm, whose conference version
appeared at STOC 2005, settles undirected connectivity~\cite{Reingold2005,Reingold2008}.
Reingold, Trevisan and Vadhan studied regular directed graphs and gave
logarithmic-space path search for Eulerian directed graphs
\cite{RTV2006}. Their oblivious walk constructions impose additional
labeling conditions. General directed configuration graphs need not
have these properties, so those results did not settle the unrestricted
randomized computation problem.
A read-once branching program is a layered deterministic computation
that reads the input symbols once in a fixed order (bits in the binary
case). Its length is the number of steps,
and its width is the largest number of states in a layer.
A $1/2$-hitting set for a class of programs contains an accepting input
for every member that accepts at least half of all input strings.
The relation between one-sided and two-sided derandomization also
requires care. Cheng and Hoza proved that one uniformly logarithmic-space-enumerable
$1/2$-hitting-set family for all width-$n$, length-$n$ read-once branching
programs would yield
$\mathsf L=\mathsf{BPL}$~\cite{ChengHoza2022}.
The premise is a hitting-set construction; it is not merely the
language-class equality $\mathsf{RL}=\mathsf L$.

Recent results refine both general and structured models.
Chen, Cohen, Doron, Khaskelberg and Ta-Shma improve error reduction for
weighted generators; in the polynomial-width regime their construction
has seed length $\log|\Sigma|+O(\log^2 T+\log(1/\delta))$, where
$T$ is branching-program length, $\Sigma$ its alphabet and $\delta$
the error~\cite{ChenCohenDoronKhaskelbergTaShma2026}.
Cheng and Wu obtain simultaneous polynomial-time and polylogarithmic-space
simulations for regular branching programs, with extensions to polynomial-time
randomized logarithmic-space machines whose two-way read-only random tape
visits each cell at most constantly often, and to such machines with read-once randomness
and an auxiliary stack whose push, pop and idle schedule is independent
of the random bits. Consecutive stays on a random-tape cell count as one
visit~\cite[Theorems~5--6]{ChengWu2026SC}.
Hardness-versus-randomness approaches additionally give conditional,
certified and per-instance alternatives; Pyne and Tell survey this
line of work~\cite{PyneTell2026}.

Effectivity has its own substantial antecedents.
Pyne, Raz and Zhan construct a universal deterministic estimator whose
$O(S(n))$ space guarantee follows from
$\mathsf{prBPL}\subseteq\mathsf{SPACE}[O(S(n))]$, for
space-constructible $S(n)\ge\log n$
\cite[full version, Theorem~5.1]{PyneRazZhan2023}.
Here $\mathsf{prBPL}$ denotes promise problems solvable in polynomial
time and randomized logarithmic space with bounded error.
The compiler in Appendix~\ref{sec:compiler} uses the fixed separator established by the preceding construction
and supplies explicit source-specific numerical resource bounds.
Its construction should be read in this context of universal
simulation, with the source bounds supplied as part of the compiler input.

